% Requires:
% \usepackage{amsmath}
% \usetikzlibrary{angles}
% \usetikzlibrary{calc}
% \usetikzlibrary{patterns}
% \usetikzlibrary{quotes}
\begin{tikzpicture}[
    scale=0.84,
    % Define a style for the angle pics based on your request
    myangle/.style={
      draw,
      angle radius=0.9cm,
      angle eccentricity=1.3,
    }
  ]

  % Parameters
  \def\angI{30}    % Incident angle
  \def\angT{20}    % Transmitted angle
  \def\rayLen{5}   % Length of rays
  \def\waveHalfLen{0.35cm} % FIXED half-length of wavefronts (total length = 0.7cm)

  % Coordinates
  \coordinate (O) at (0,0);

  % Ray Endpoints
  % Incident comes from bottom-left (180 + 30 = 210 degrees)
  \coordinate (IncStart) at ({180+\angI}:\rayLen);
  % Reflected goes to top-left (180 - 30 = 150 degrees)
  \coordinate (RefEnd) at ({180-\angI}:\rayLen);
  % Transmitted goes to top-right (20 degrees)
  \coordinate (TransEnd) at (\angT:\rayLen);

  % Normal reference points for defining angles
  \coordinate (N_left) at (-4,0);
  \coordinate (N_right) at (4,0);

  % Mediums
  \node[align=center] at (-2.5, 3.3) {Medium 1\\$\mu_1, \varepsilon_1$};

  \node[align=center] at (2.5, 3.3) {Medium 2\\$\mu_2, \varepsilon_2$};

  % Interface and Normal
  % Hatching on the right side
  \fill[pattern=north east lines] (0,-4) rectangle (0.3, 4);
  \draw[thick] (0,-4) -- (0,4);

  % Normal
  \draw[dashed] (-5,0) -- (5,0) node[pos=.75,below=5pt] {Surface normal};

  % Rays
  % Incident
  \draw[thick, ->] (IncStart) -- (O) node[pos=.65,align=center, below,sloped,font=\footnotesize] {Incident\\wave};

  % Reflected
  \draw[thick, ->] (O) -- (RefEnd) node[pos=.35,align=center, above,sloped,font=\footnotesize] {Reflected\\wave};

  % Transmitted
  \draw[thick, ->] (O) -- (TransEnd) node[pos=.35,align=center, above,sloped,font=\footnotesize] {Transmitted\\(refracted) wave};

  % Wavefronts
  % Using ($ (Position)!length!angle:(Target) $) syntax for fixed orthogonal length

  % Incident Wavefronts
  \foreach \t in {0.10, 0.13, 0.16} {
    \coordinate (W) at ($(IncStart)!\t!(O)$);
    \draw[dashed,thick,gray] ($(W)!\waveHalfLen!90:(O)$) -- ($(W)!\waveHalfLen!-90:(O)$);
  }

  % Reflected Wavefronts
  \foreach \t in {0.64, 0.67, 0.70} {
    \coordinate (W) at ($(O)!\t!(RefEnd)$);
    \draw[dashed,thick,gray] ($(W)!\waveHalfLen!90:(RefEnd)$) -- ($(W)!\waveHalfLen!-90:(RefEnd)$);
  }

  % Transmitted Wavefronts
  \foreach \t in {0.64, 0.67, 0.70} {
    \coordinate (W) at ($(O)!\t!(TransEnd)$);
    \draw[dashed,thick,gray] ($(W)!\waveHalfLen!90:(TransEnd)$) -- ($(W)!\waveHalfLen!-90:(TransEnd)$);
  }

  % Angles

  % Incident Angle (between Incident Ray and Normal)
  % Note: defined by points IncStart--O--N_left
  \draw pic [|->,myangle, "$\theta_i$"] {angle = N_left--O--IncStart};

  % Reflected Angle (between Normal and Reflected Ray)
  % Note: defined by points N_left--O--RefEnd
  \draw pic [<-|,myangle, "$\theta_r$"] {angle = RefEnd--O--N_left};

  % Transmitted Angle (between Normal and Transmitted Ray)
  % Note: defined by points TransEnd--O--N_right
  % Adjusted radius slightly to avoid overlapping the wavefronts
  \draw pic [->,myangle, "$\theta_t$", angle radius=1.1cm] {angle = N_right--O--TransEnd};

  % Electric Field Vectors

  % Incident E-field
  \coordinate (EI_pos) at ($(IncStart)!0.00!(O)$);
  % E_parallel (Arrow perp to ray)
  \draw[->,thick,blue] (EI_pos) -- ($(EI_pos)!0.7cm!90:(O)$) node[above] {$E_\parallel^i$};
  % E_perp (Circle with dot)
  \draw[red,fill=white] (EI_pos) circle (6pt);
  \fill[red] (EI_pos) circle (2pt);
  \node[red,below=5pt] at (EI_pos) {$E_\perp^i$};

  % Reflected E-field
  \coordinate (ER_pos) at ($(O)!0.8!(RefEnd)$);
  % E_parallel (Arrow perp to ray)
  \draw[->,thick,blue] (ER_pos) -- ($(ER_pos)!0.7cm!+90:(RefEnd)$) node[left] {$E_\parallel^r$};
  % E_perp
  \draw[red,fill=white] (ER_pos) circle (6pt);
  \fill[red] (ER_pos) circle (2pt);
  \node[red,above=5pt] at (ER_pos) {$E_\perp^r$};

  % Transmitted E-field
  \coordinate (ET_pos) at ($(O)!0.8!(TransEnd)$);
  % E_parallel
  \draw[->,thick,blue] (ET_pos) -- ($(ET_pos)!0.7cm!90:(TransEnd)$) node[above] {$E_\parallel^t$};
  % E_perp
  \draw[red,fill=white] (ET_pos) circle (6pt);
  \fill[red] (ET_pos) circle (2pt);
  \node[red,below=8pt,xshift=3pt] at (ET_pos) {$E_\perp^t$};

  % Legend
  \begin{scope}[shift={(1.2, -3.5)}]
    \draw[dashed,thick,draw=gray] (0,0.4) -- ++ ({2*\waveHalfLen},0) node[right] {Wavefronts};
  \end{scope}

\end{tikzpicture}
